Table~\ref{tab:performance} gives the repeated comparison, and Fig.~\ref{fig:latency} shows the variation in GET and POST P95 separately. The 600 requests/s comparison met the specified criterion in all ten trials per policy. Its median in-window delivery was at least 99.9\% for every policy.

At 4{,}800 offered requests/s, B had the lowest median item GET P95, 4.14\,ms, and the highest median goodput, 4{,}777.3 successful requests/s. Its median POST P95 was 54.30\,ms, compared with 35.31\,ms in A. Table~\ref{tab:paired} shows the paired changes. The POST difference for B relative to A had the opposite sign from the GET difference. The load criterion uses overall P95, so an individual endpoint can exceed 50\,ms even when a trial passes. A passed nine of ten trials, B and C passed ten each, and D passed seven. The recorded load shortfalls came from scheduled starts that k6 dropped and responses that finished outside the measurement window. The repeated open-loop trials recorded no failed or incomplete HTTP attempts after a start.

At 4{,}800 requests/s, D used a median 524.7 item SELECTs per 1{,}000 started GETs, compared with 424.4 in C and 1{,}000 in A. D counted a median of 3 rejected generation checks per trial at 600 requests/s and 3{,}417.5 at 4{,}800 requests/s. A generation mismatch shows a refill crossed an invalidation. It can conservatively reject a candidate that has already seen the newer commit. The paired median change in delivery from C to D was about \(-1.6\) percentage points. The paired GET and POST P95 intervals for that comparison spanned zero. The data support a delivery cost for D within this test setup. They leave the separate effect of Redis command work on latency unresolved. Omitting repeat 1 preserved the ordering of the four delivery medians.

In the TTL sweep (Fig.~\ref{fig:ttl}), B's median payload hit rate increased from 86.7\% at 1\,s to 98.1\% at 10\,s. Its median GET P95 fell from 1.04 to 0.74\,ms. D's median hit rate stayed between 62.4\% and 63.8\%, with median GET P95 near 1.2\,ms at all three TTLs. That weaker TTL sensitivity is consistent with orders removing payloads before expiry. The sweep measured hit rate and response time at 600 requests/s.
